% Part: proof-theory
% Chapter: sequent-calculus
% Section: rules-G1c

\documentclass[../../../include/open-logic-section]{subfiles}

\begin{document}

\begin{table}
\begin{tabular}{@{}c@{}c@{}}
\multicolumn{2}{@{}c@{}}{स्वयंसिद्धके}\\[2ex]
$!A \Sequent !A$ & $\lfalse \Sequent \quad$
\\[2ex]
\multicolumn{2}{@{}c@{}}{संरचनात्मक नियम}\\[2ex]
\Axiom$ \Gamma \fCenter \Delta $
\RightLabel{\LeftR{\Weakening}}
\UnaryInf$ !A, \Gamma \fCenter \Delta$
\DisplayProof
&
\Axiom$ \Gamma \fCenter $
\RightLabel{\RightR{\Weakening}}
\UnaryInf$ \Gamma \fCenter !A$
\DisplayProof
\\[4ex]
\Axiom$ !A, !A, \Gamma \fCenter \Delta $
\RightLabel{\LeftR{\Contraction}}
\UnaryInf$ !A, \Gamma \fCenter \Delta$
\DisplayProof
&
\\[4ex]
\multicolumn{2}{@{}c@{}}{तार्किक नियम}\\[2ex]
\Axiom$ \Gamma \fCenter !A $
\RightLabel{\LeftR{\lnot}}
\UnaryInf$ \lnot !A, \Gamma \fCenter $
\DisplayProof
&
\Axiom$!A, \Gamma \fCenter $
\RightLabel{\RightR{\lnot}}
\UnaryInf$ \Gamma \fCenter \lnot !A $
\DisplayProof
\\[4ex]\begin{tabular}[t]{@{}l@{}}
\Axiom$ !A, \Gamma \fCenter \Delta$
\RightLabel{\LeftR{\land}}
\UnaryInf$ !A \land !B, \Gamma \fCenter \Delta$
\DisplayProof
\\[3ex]
\Axiom$!B, \Gamma \fCenter \Delta$
\RightLabel{\LeftR{\land}}
\UnaryInf$!A \land !B, \Gamma \fCenter \Delta$
\DisplayProof\\[3ex]
\end{tabular}
&
\Axiom$\Gamma \fCenter !A$
\Axiom$ \Gamma \fCenter !B$
\RightLabel{\RightR{\land}}
\BinaryInf$ \Gamma \fCenter !A \land !B $
\DisplayProof
\\[4ex]\Axiom$!A, \Gamma \fCenter \Delta$
\Axiom$!B, \Gamma \fCenter \Delta$
\RightLabel{\LeftR{\lor}}
\BinaryInf$!A \lor !B, \Gamma \fCenter \Delta$
\DisplayProof
&
\begin{tabular}[t]{@{}r@{}}
\Axiom$\Gamma \fCenter !A$
\RightLabel{\RightR{\lor}}
\UnaryInf$ \Gamma \fCenter !A \lor !B$
\DisplayProof
\\[3ex]
\Axiom$ \Gamma \fCenter !B$
\RightLabel{\RightR{\lor}}
\UnaryInf$ \Gamma \fCenter !A \lor !B$
\DisplayProof\\[3ex]
\end{tabular}
\\[4ex]
\Axiom$ \Gamma \fCenter !A$
\Axiom$ !B, \Gamma \fCenter \Delta$
\RightLabel{\LeftR{\lif}}
\BinaryInf$ !A \lif !B, \Gamma \fCenter \Delta$
\DisplayProof
&
\Axiom$ !A, \Gamma \fCenter !B$
\RightLabel{\RightR{\lif}}
\UnaryInf$ \Gamma \fCenter !A \lif !B $
\DisplayProof
\\[4ex]
\Axiom$ !A(t), \Gamma \fCenter \Delta$
\RightLabel{\LeftR{\lforall}}
\UnaryInf$ \lforall[x][!A(x)],\Gamma \fCenter \Delta$
\DisplayProof
&
\Axiom$ \Gamma \fCenter !A(a) $
\RightLabel{\RightR{\lforall}}
\UnaryInf$ \Gamma \fCenter \lforall[x][!A(x)]$
\DisplayProof
\\[4ex]
\Axiom$ !A(a), \Gamma \fCenter \Delta $
\RightLabel{\LeftR{\lexists}}
\UnaryInf$ \lexists[x][!A(x)], \Gamma \fCenter \Delta$
\DisplayProof
&
\Axiom$ \Gamma \fCenter !A(t) $
\RightLabel{\RightR{\lexists}}
\UnaryInf$ \Gamma \fCenter \lexists[x][!A(x)]$
\DisplayProof
\end{tabular}
\caption{अंतःप्रज्ञावादी तर्कशास्त्रासाठी \Log{G1i} चे नियम.
$\RightR{\forall}$ आणि $\LeftR{\exists}$ मध्ये !!{constant}~$a$ हा
निष्कर्ष-क्रमवर्तीमध्ये येऊ नये. $\Gamma$ हा !!{formula}s चा
बहुसंच असून $\Delta$ रिकामा असू शकतो किंवा त्यात एकच
!!{formula} असू शकते. \RightR{\Contraction} नियम नाही. \Log{G1m}
म्हणजे~\RightR{\Weakening} आणि $\lfalse,
\Gamma \Sequent \Delta$ या रूपाची स्वयंसिद्धके वगळलेली \Log{G1i}.}
\ollabel{tab:G1i}
\end{table}

\end{document}
